
\SGARefTarget{sga3:xxiv:1.3}
\paragraph{\SGARefLink{sga3:xxiv:1.3}{Theorem 1.3}.}
\emph{Let \(S\) be a scheme and \(G\) a reductive \(S\)-group. Consider
the canonical exact sequence of \(S\)-sheaves of groups}
\footnote{Editor's note: Recall that
\(\operatorname{ad}(G)=G/\underline{\Centr}(G)\) denotes the adjoint
group of \(G\); cf.\ \SGARefLink{sga3:xxii:def:4.3.6}{Exp.~XXII, 4.3.6}.}
\[
  1\longrightarrow \operatorname{ad}(G)
  \longrightarrow \underline{\Aut}_{S\text{-}\mathrm{gr.}}(G)
  \xrightarrow{p}\underline{\Autext}(G)
  \longrightarrow 1.
\]
\emph{Then:}
\begin{enumerate}
\item[\SGAInlineRefTarget{sga3:xxiv:1.3:item:i}(i)] \(\underline{\Aut}_{S\text{-}\mathrm{gr.}}(G)\) is
  representable by a smooth separated \(S\)-scheme;
\item[\SGAInlineRefTarget{sga3:xxiv:1.3:item:ii}(ii)] \(\underline{\Autext}(G)\) is representable by a twisted
  constant \(S\)-scheme of finite generation \textup{(\SGARefLink{sga3:x:definition:5-1}{Exp.~X, 5.1})};
\item[\SGAInlineRefTarget{sga3:xxiv:1.3:item:iii}(iii)] if \(G\) is split, the preceding exact sequence is split.
  More precisely, for every pinning of \(G\), the morphism
  \textup{(cf.\ \SGARefLink{sga3:xxiv:1.2}{1.2}~\((*)\))}
  \[
    p\circ a:E_S\longrightarrow\underline{\Autext}(G)
  \]
  is an isomorphism.
\end{enumerate}

Let us first show how the theorem follows from the following lemma.

\SGARefTarget{sga3:xxiv:1.4}
\paragraph{\SGARefLink{sga3:xxiv:1.4}{Lemma 1.4}.}
\emph{Under the hypotheses of \textup{\SGARefLink{sga3:xxiv:1.3:item:iii}{(iii)}},
\(\underline{\Aut}_{S\text{-}\mathrm{gr.}}(G)\) is the semidirect
product}
\[
  a(E_S)\cdot\operatorname{ad}(G).
\]
\footnote{Editor's note: Here and below, the notation
\(\operatorname{int}(G)\) has been replaced by
\(\operatorname{ad}(G)\).}

The lemma immediately implies the theorem when \(G\) is split. Since
\(G\) is locally split for the étale topology \textup{(\SGARefLink{sga3:xxii:cor:2.3}{Exp.~XXII,
2.3})}, hence also for the \textup{(fppf)} topology, and since the
latter is ``of effective descent'' for the fibred category of twisted
constant morphisms \textup{(\SGARefLink{sga3:x:corollary:5-5}{Exp.~X, 5.5})}, assertion \textup{\SGARefLink{sga3:xxiv:1.3:item:ii}{(ii)}}
follows in the general case \textup{(cf.\ \SGARefLink{sga3:iv:proposition:4:6:8}{Exp.~IV, 4.6.8})}. To deduce
\textup{\SGARefLink{sga3:xxiv:1.3:item:i}{(i)}}, observe that \(\operatorname{ad}(G)\) is affine over
\(S\); thus \(p\) is affine whenever
\(\underline{\Aut}_{S\text{-}\mathrm{gr.}}(G)\) is representable, and
the result follows by descent for affine schemes.
\footnote{Editor's note: Cf.\ SGA~1, VIII, \SGARefLink{sga3:viii:proposition:2:1}{2.1}.}

It remains only to prove \SGARefLink{sga3:xxiv:1.4}{1.4}. For this, it suffices to prove:

\SGARefTarget{sga3:xxiv:1.5}
\paragraph{\SGARefLink{sga3:xxiv:1.5}{Lemma 1.5}.}
\emph{If \((\mathscr R,\mathcal E)\) and
\((\mathscr R',\mathcal E')\) are two pinnings of the reductive
\(S\)-group \(G\), there exists a unique inner automorphism \(u\) of
\(G\) over \(S\) carrying one pinning to the other; that is, such that
there exists an isomorphism
\[
  v:\mathscr R'\xrightarrow{\sim}\mathscr R
\]
for which \(u(\mathcal E)=v(\mathcal E')\); cf.\ \SGARefLink{sga3:xxiv:1.0}{1.0}.}

\SGARefTarget{sga3:xxiv:1.5.1}
\subsubsection*{1.5.1. Uniqueness.}
It suffices to prove that, if \(G\) is a pinned \(S\)-group and
\(\operatorname{int}(g)\) is an automorphism of the pinned group
\textup{(\(g\in G(S)\))}, then \(\operatorname{int}(g)=\operatorname{id}\).
First,
\[
  \operatorname{int}(g)T=T,\qquad
  \operatorname{int}(g)B=B,
\]
and hence
\[
  g\in
  \underline{\operatorname{Norm}}_G(T)(S)
  \cap\underline{\operatorname{Norm}}_G(B)(S)
  =T(S)
\]
\textup{(cf., for example, \SGARefLink{sga3:xxii:prop:5.6.1}{Exp.~XXII, 5.6.1})}. It follows that
\(\operatorname{int}(g)\) normalizes every \(U_\alpha\) and that
\[
  \Lie(\operatorname{int}(g))X_\alpha
  =\operatorname{Ad}(g)X_\alpha
  =\alpha(g)X_\alpha
\]
for every \(\alpha\in\Delta\). Thus \(\alpha(g)=1\) for
\(\alpha\in\Delta\), and therefore
\[
  g\in
  \bigcap_{\alpha\in\Delta}(\ker\alpha)(S)
  =\underline{\Centr}(G)(S)
\]
\textup{(\SGARefLink{sga3:xxii:rem:4.1.8}{Exp.~XXII, 4.1.8})}. \qed
